%--- iliad ---
% summary: >-
%   How an idealized agent should predict. A Bayesian mixture over a countable
%   class of computable hypotheses learns to predict any sequence, with total error
%   bounded by the description length of the truth -- a formal Occam's razor.
%--- end ---
\documentclass[fontsize=11pt, parskip=half]{scrartcl}
% Page geometry (compact screen page, 1cm margins) now comes from iliad.sty.
\pagestyle{plain}

\usepackage{amsmath, amssymb, amsthm, amsfonts}

\usepackage[dvipsnames]{xcolor}


% iliad.sty: local copy first (Overleaf), else the shared ../iliad.sty
\IfFileExists{iliad.sty}{\usepackage[boxes]{iliad}}{\usepackage[boxes]{../iliad}}

% sheet notation
\renewcommand{\thesection}{\arabic{section}}
\newcommand{\Msol}{\mathcal{M}_{sol}}
\newcommand{\Mc}{\mathcal{M}}
\newcommand{\KL}{D_{\text{KL}}}
\newcommand{\Sn}[2]{S_{n}(#1 \parallel #2)}

\usepackage{cancel}

% 3. Define your theorem environments AFTER cleveref is loaded
%===================================
\begin{document}

\title{Solomonoff Induction}
\author{\authorname{David Quarel}\\ \affiliation{Australian National University}
  \and \authorname{Leon Lang}\\ \affiliation{University of Amsterdam}}
\date{\today}
\maketitle

\begin{learningoutcomes}
  \begin{itemize}
    \item Define the Bayesian mixture over a class of environments and prove it is a proper predictor whose posterior updates multiplicatively.
    \item Prove the cumulative prediction-error bound and specialize it to the Solomonoff prior to recover the Occam bound.
    \item Show the mixture makes boundedly many mistakes and is Pareto-optimal under KL and squared loss.
    \item Bound the misspecified case when the truth lies outside the model class.
  \end{itemize}
\end{learningoutcomes}

\medskip
\noindent
Much of this material is drawn from \cite[Chapter 3]{Hutter:24uaibook2}
and the earlier \cite{Hutter:04uaibook}; both books provide fuller explanations, additional context, and proofs that we gloss over here.

\medskip
\noindent
\textbf{Difficulty ratings.} Each problem is tagged with a rating in square brackets following Knuth's exercise rating scheme~\cite{Knuth:73a}: roughly, \textbf{[00]}~trivial, \textbf{[10]}~15--minute pencil-and-paper, \textbf{[20]}~1--2~hours, \textbf{[30]}~several hours to a day, \textbf{[40]}~a significant research result. Intermediate values are possible.

%===================================

\section*{Overview}

Solomonoff induction is the \emph{prediction} half of \textbf{Universal AI}: what
would an \emph{optimally intelligent} agent do if it only had to predict, given
unlimited compute and the weakest possible assumptions? The answer is a single
\textbf{Bayesian mixture} $\xi$ over a countable class $\mathcal{M}$ of candidate
environments, weighted by a prior $w_\nu$. It (eventually) predicts as well as the
true environment $\mu$, and under the \textbf{universal} choice --- $\mathcal{M}$ =
all computable environments, prior $w_\nu = 2^{-K(\nu)}$ with $K$ the Kolmogorov
complexity --- the assumption ``$\mu \in \mathcal{M}$'' becomes ``the universe is
computable,'' and Occam's razor drops out of the mathematics.

This module is a worksheet: you build the theory yourself, one problem at a time.
The sequel module, \href{/agency/aixi}{AIXI}, lifts the same mixture to sequential
decision-making (learning to \emph{act}).

\section*{Prerequisites}

\begin{itemize}
    \item Comfort with discrete probability: conditional distributions, the chain
    rule, expectations.
    \item Kullback--Leibler divergence and basic information theory (helpful,
    developed as needed).
\end{itemize}

\section*{Goal and Roadmap}

\textbf{A recipe for prediction:} Solomonoff induction is an attempt to mathematically
formalize the \href{https://en.wikipedia.org/wiki/Problem_of_induction}{\emph{problem of induction}}
from philosophy: How to make predictions about the future based on past observations?

We formalize this as sequence prediction: There is some true environment $\mu$ that generates
a sequence $x_1, x_2, \dots$ of (binary) symbols. Each next symbol $x_t \sim \mu(\cdot \mid x_{<t})$ is sampled
from $\mu$ conditioned on the past $x_{<t}$. A predictor $P$ takes the history $x_{<t}$ and gives a 
distribution $P(\cdot \mid x_{<t})$ over the next symbol $x_t$. 

We measure the quality of a predictor $P$ by the $\mu$-expected squared prediction error, summed over every timestep:
\begin{align}
S_\infty^\mu := \sum_{t=1}^{\infty} \sum_{x_{<t} \in \mathbb{B}^*} \mu(x_{<t}) \sum_{x_t \in \mathbb{B}} \big( P(x_t \mid x_{<t}) - \mu(x_t \mid x_{<t}) \big)^2.
\end{align}

How to construct a predictor $P$ such that $S^\mu_\infty$ is small (or at least finite)? Bayesian inference to the rescue: we choose as our predictor a \emph{Bayesian mixture} $\xi$ over a countable class $\Mc = \{\nu_1, \nu_2, \dots\}$ of candidate environments, weighted by prior beliefs $w_\nu$ (formalized in \cref{def:mixture}). The main results we build up to are:
\begin{itemize}\itemsep=2pt
  \item \textbf{Cumulative bound} (\cref{prob:bound}): Assuming $\mu \in \Mc$, $S^\mu_\infty \leq -\ln w_\mu$. The higher the prior
  on $\mu$, the lower the prediction error.
  \item \textbf{Explicit bound} (\cref{prob:bound_explicit_ex}): specialize to the Solomonoff prior $w_\nu = 2^{-K(\nu)}$ to get $S^\mu_\infty \leq K(\mu)\ln 2$, where $K$ is the \emph{Kolmogorov complexity}.
  \item \textbf{Pareto optimality} (\cref{prob:kl_pareto,prob:sq_pareto}): no other predictor weakly dominates $\xi$ on every $\nu \in \Mc$, for either KL or squared loss.
  \item \textbf{Misspecified version} (\cref{prob:misspec}): If $\mu \not\in \Mc$, the cumulative bound becomes $-\ln w_{\hat\mu} + D_n(\mu \parallel \hat\mu)$: the constant complexity term plus an approximation term $D_n(\mu \parallel \hat\mu) = \sum_{t=1}^n d_t(\mu \parallel \hat\mu)$ that in general grows linearly in $n$, so $S^\mu_\infty$ diverges. Here $\hat{\mu} \in \Mc$ is the ``closest'' environment to $\mu$.
\end{itemize}

%===================================
\section*{Notation}

For more background, see \href{https://www.lesswrong.com/posts/HSDumToH57nSRdLST/a-technical-introduction-to-solomonoff-induction-without-k}{the corresponding post on Solomonoff induction}.

\begin{callout}[note]
\textbf{Symbols used throughout.}
\begin{itemize}\itemsep=2pt
  \item $\mathbb{B} = \{\texttt{0}, \texttt{1}\}$: the binary alphabet; $\mathbb{B}^*$ all finite binary strings; $\mathbb{B}^n$ length-$n$ strings.
  \item $\epsilon$: the empty string. $x_{1:n} := x_1 x_2 \cdots x_n$; $x_{<t} := x_1 x_2 \cdots x_{t-1}$.
  \item $X_{1:n}, X_{<t}$: the corresponding random variables.
  \item $\nu, \rho$: generic environments / predictors. $\mu$: the true (unknown) environment generating the data. $\xi$: a Bayesian mixture of environments.
  \item $\Mc = \{\nu_1, \nu_2, \ldots\}$: a countable class of candidate environments.
  \item $w_\nu > 0$ with $\sum_{\nu \in \Mc} w_\nu = 1$: the prior over $\Mc$. $w(\nu \mid x_{<t})$: the posterior after observing $x_{<t}$.
  \item $\Delta \mathcal{X}$: the set of probability distributions over a finite set $\mathcal{X}$.
\end{itemize}
\end{callout}

%===================================
\section{The mixture is a predictor}\label{prob:one_step}

\begin{definition}[Environment]\label{def:environment}
An \textbf{environment} $\nu$ assigns to each history $x_{<t} \in \mathbb{B}^*$ a predictive distribution $\nu(\cdot \mid x_{<t}) \in \Delta \mathbb{B}$ over the next symbol. The joint probability of a string is defined by the chain rule
\[
\nu(x_{1:n}) ~:=~ \prod_{t=1}^{n} \nu(x_t \mid x_{<t}), \qquad \nu(\epsilon) := 1,
\]
so that $\nu(x_{1:t}) = \nu(x_{<t}) \cdot \nu(x_t \mid x_{<t})$ and $\nu(x) = \nu(x\texttt{0}) + \nu(x\texttt{1})$ for all $x \in \mathbb{B}^*$. The true (unknown) environment generating the data is denoted $\mu$.
\end{definition}

\begin{definition}[Bayesian mixture $\xi$]\label{def:mixture}
Let $\Mc = \{\nu_1, \nu_2, \dots\}$ be a countable class of environments with prior weights $w_\nu > 0$ satisfying $\sum_{\nu \in \Mc} w_\nu = 1$. The \textbf{Bayesian mixture} is defined as a prior-weighted mixture over all environments $\nu$ in $\Mc$:
\begin{equation}\label{eq:mixture_expression}
  \xi(x_{1:t}) ~:=~ \sum_{\nu \in \Mc} w_\nu\, \nu(x_{1:t}).
\end{equation}
Its one-step predictive distribution is
\[
  \xi(x_t \mid x_{<t}) ~=~ \sum_{\nu \in \Mc} w(\nu \mid x_{<t})\, \nu(x_t \mid x_{<t}),
\]
with posterior weights 
\[
w(\nu \mid x_{<t}) ~:=~ w_\nu\, \frac{\nu(x_{<t})}{\xi(x_{<t})}
\qquad w(\nu \mid \epsilon) := w_\nu.
\]
Here $w(\nu \mid x_{<t})$ is the \emph{posterior} belief in $\nu$ after observing $x_{<t}$.
\end{definition}

The mixture $\xi$ is defined as a sum of \emph{joint} probabilities. To use it as a predictor we need its one-step conditional $\xi(x_t \mid x_{<t})$, and we need to know it is a genuine probability distribution (so that, later, Pinsker's inequality applies to it).

\begin{exercise}[Generalized chain rule]
\label{prob:chain_rule_gen_ex}
\difficulty{05}



The two-term chain rule baked into \cref{def:environment} extends to arbitrary contiguous blocks. Show that for any $1 \leq p \leq q \leq r$ and any history $x_{<p}$ with $\nu(x_{<p}) > 0$,
\[
  \nu(x_{p:r} \mid x_{<p}) ~=~ \nu(x_{p:q} \mid x_{<p}) \cdot \nu(x_{q+1:r} \mid x_{<p}\, x_{p:q}).
\]
\end{exercise}

\begin{solution}[prob:chain_rule_gen_ex]
Expand the conditional as a ratio of joints; the chain rule (\cref{def:environment}) telescopes the ratio to a product over indices $k = p, \ldots, r$:
\[
  \nu(x_{p:r} \mid x_{<p})
  ~=~ \frac{\nu(x_{1:r})}{\nu(x_{<p})}
  ~=~ \frac{\prod_{k=1}^{r} \nu(x_k \mid x_{<k})}{\prod_{k=1}^{p-1} \nu(x_k \mid x_{<k})}
  ~=~ \prod_{k=p}^{r} \nu(x_k \mid x_{<k}).
\]
For any $k \geq p$, the past $x_{<k}$ is just $x_{<p}$ extended by $x_{p:k-1}$, so the same expansion identifies every contiguous sub-block as a conditional with the right history. Splitting the product at index $q$:
\[
  \nu(x_{p:r} \mid x_{<p})
  ~=~ \underbrace{\prod_{k=p}^{q} \nu(x_k \mid x_{<k})}_{=\,\nu(x_{p:q}\,\mid\,x_{<p})}
  \;\cdot\; \underbrace{\prod_{k=q+1}^{r} \nu(x_k \mid x_{<k})}_{=\,\nu(x_{q+1:r}\,\mid\,x_{<p}\, x_{p:q})}. \qedhere
\]
\end{solution}

\begin{exercise}[Mixture is a predictor]
\label{prob:one_step_ex}
\difficulty{10}



Starting from $\xi(x_{1:t}) = \sum_{\nu \in \Mc} w_\nu\, \nu(x_{1:t})$, show that
\[
\xi(x_t \mid x_{<t}) ~=~ \sum_{\nu \in \Mc} w(\nu \mid x_{<t})\, \nu(x_t \mid x_{<t}),
\qquad w(\nu \mid x_{<t}) := w_\nu\, \frac{\nu(x_{<t})}{\xi(x_{<t})},
\]
and conclude that $\sum_{x_t \in \mathbb{B}} \xi(x_t \mid x_{<t}) = 1$, i.e.\ $\xi(\cdot \mid x_{<t})$ is a probability distribution over $\mathbb{B}$.

\begin{hint}
Write $\xi(x_t \mid x_{<t}) = \xi(x_{1:t}) / \xi(x_{<t})$, expand the numerator, and use the chain rule $\nu(x_{1:t}) = \nu(x_{<t})\, \nu(x_t \mid x_{<t})$ (\cref{def:environment}). For the last part, note that the posterior weights sum to $1$.
\end{hint}
\end{exercise}

\begin{solution}[prob:one_step_ex]
The one-step conditional is the ratio of joint to marginal, $\xi(x_t \mid x_{<t}) := \xi(x_{1:t}) / \xi(x_{<t})$. Expanding the numerator and applying the chain rule $\nu(x_{1:t}) = \nu(x_{<t})\, \nu(x_t \mid x_{<t})$ to each term:
\begin{align*}
\xi(x_t \mid x_{<t})
~&=~ \frac{\sum_{\nu \in \Mc} w_\nu\, \nu(x_{1:t})}{\xi(x_{<t})}
~=~ \frac{\sum_{\nu \in \Mc} w_\nu\, \nu(x_{<t})\, \nu(x_t \mid x_{<t})}{\xi(x_{<t})} \\
~&=~ \sum_{\nu \in \Mc} \underbrace{\frac{w_\nu\, \nu(x_{<t})}{\xi(x_{<t})}}_{=\, w(\nu \mid x_{<t})}\, \nu(x_t \mid x_{<t})
~=~ \sum_{\nu \in \Mc} w(\nu \mid x_{<t})\, \nu(x_t \mid x_{<t}).
\end{align*}
The posterior weights are non-negative and sum to $1$:
\[
\sum_{\nu \in \Mc} w(\nu \mid x_{<t})
~=~ \frac{\sum_{\nu \in \Mc} w_\nu\, \nu(x_{<t})}{\xi(x_{<t})}
~=~ \frac{\xi(x_{<t})}{\xi(x_{<t})}
~=~ 1.
\]
Therefore $\xi(\cdot \mid x_{<t})$ is a convex combination of the probability distributions $\nu(\cdot \mid x_{<t})$, so it is itself a probability distribution over $\mathbb{B}$:
\[
\sum_{x_t \in \mathbb{B}} \xi(x_t \mid x_{<t})
~=~ \sum_{\nu \in \Mc} w(\nu \mid x_{<t}) \underbrace{\sum_{x_t \in \mathbb{B}} \nu(x_t \mid x_{<t})}_{=\, 1}
~=~ \sum_{\nu \in \Mc} w(\nu \mid x_{<t})
~=~ 1. \qedhere
\]
\end{solution}

\begin{exercise}[Posterior update]
\label{prob:posterior_update_ex}
\difficulty{10}



Show that the posterior weight $w(\cdot \mid \cdot)$ updates \emph{multiplicatively} as new symbols arrive. For every $x_{1:t} \in \mathbb{B}^*$ with $\xi(x_{1:t}) > 0$,
\[
  w(\nu \mid x_{1:t}) ~=~ w(\nu \mid x_{<t}) \cdot \frac{\nu(x_t \mid x_{<t})}{\xi(x_t \mid x_{<t})}.
\]
That is: the new posterior equals the old posterior, scaled by the \textbf{likelihood ratio} of how well $\nu$ predicted the just-observed symbol relative to the mixture's own prediction.
\end{exercise}

\begin{solution}[prob:posterior_update_ex]
By \cref{def:mixture},
\[
  w(\nu \mid x_{1:t}) ~=~ w_\nu \,\frac{\nu(x_{1:t})}{\xi(x_{1:t})},
  \qquad
  w(\nu \mid x_{<t}) ~=~ w_\nu \,\frac{\nu(x_{<t})}{\xi(x_{<t})}.
\]
Apply the chain rule (\cref{def:environment}) to both $\nu$ and $\xi$: $\nu(x_{1:t}) = \nu(x_{<t})\,\nu(x_t \mid x_{<t})$ and $\xi(x_{1:t}) = \xi(x_{<t})\,\xi(x_t \mid x_{<t})$. Dividing,
\[
  \frac{w(\nu \mid x_{1:t})}{w(\nu \mid x_{<t})}
  ~=~ \frac{\nu(x_{1:t})}{\nu(x_{<t})} \cdot \frac{\xi(x_{<t})}{\xi(x_{1:t})}
  ~=~ \frac{\nu(x_t \mid x_{<t})}{\xi(x_t \mid x_{<t})}.
\]
Rearranging gives the claim. \qedhere
\end{solution}

\begin{exercise}[Multi-step posterior linearity]
\label{prob:multi_step_xi_ex}
\difficulty{10}



\Cref{prob:one_step_ex} showed that $\xi$'s one-step prediction $\xi(x_t \mid x_{<t})$ is the posterior-weighted average of the per-model one-step predictions. The same identity extends to predictions over a \emph{whole future segment} $x_{t:m}$: for any $t \leq m$,
\[
  \xi(x_{t:m} \mid x_{<t}) ~=~ \sum_{\nu \in \Mc} w(\nu \mid x_{<t})\, \nu(x_{t:m} \mid x_{<t}).
\]
That is: the \emph{same} posterior $w(\nu \mid x_{<t})$ governs $\xi$'s predictions for arbitrarily many steps into the future, not just the next one. Show this.

\begin{hint}
Write $\xi(x_{t:m} \mid x_{<t}) = \xi(x_{1:m}) / \xi(x_{<t})$, expand the numerator using the mixture definition, and apply the generalized chain rule (\cref{prob:chain_rule_gen_ex} with $p=1$, $q=t-1$, $r=m$) to factor each $\nu(x_{1:m}) = \nu(x_{<t})\,\nu(x_{t:m} \mid x_{<t})$.
\end{hint}
\end{exercise}

\begin{solution}[prob:multi_step_xi_ex]
By definition of the conditional and the mixture form $\xi(x_{1:m}) = \sum_\nu w_\nu \nu(x_{1:m})$ (\cref{def:mixture}),
\[
  \xi(x_{t:m} \mid x_{<t})
  ~=~ \frac{\xi(x_{1:m})}{\xi(x_{<t})}
  ~=~ \frac{\sum_\nu w_\nu\, \nu(x_{1:m})}{\xi(x_{<t})}.
\]
The generalized chain rule (\cref{prob:chain_rule_gen_ex} with $p=1$, $q=t-1$, $r=m$) gives $\nu(x_{1:m}) = \nu(x_{<t})\,\nu(x_{t:m} \mid x_{<t})$. Substituting,
\begin{align*}
  \xi(x_{t:m} \mid x_{<t})
  ~&=~ \sum_\nu \underbrace{\frac{w_\nu\, \nu(x_{<t})}{\xi(x_{<t})}}_{=\, w(\nu \mid x_{<t})}\, \nu(x_{t:m} \mid x_{<t}) \\
  ~&=~ \sum_{\nu \in \Mc} w(\nu \mid x_{<t})\, \nu(x_{t:m} \mid x_{<t}). \qedhere
\end{align*}
\end{solution}

%===================================
\section{KL divergence}\label{prob:chain_rule}

% Formal KL definitions are commented out: the worksheet only ever uses the
% cumulative/per-step forms D_n and d_t (\cref{def:Dn}) below, never these
% raw single- and two-variable forms. The two-variable chain rule for KL --
% formerly its own [10] exercise -- is now folded inline into the induction
% step of \cref{prob:telescope_ex}.
%
% \begin{definition}[Kullback--Leibler divergence]\label{def:kl}
% Let $P$ and $Q$ be probability distributions over a finite set $\mathcal{X}$. The \textbf{Kullback--Leibler (KL) divergence} from $P$ to $Q$ is
% \[
% D_{\text{KL}}(P \,\|\, Q) := \sum_{x \in \mathcal{X}} P(x) \ln \frac{P(x)}{Q(x)},
% \]
% where we define $0 \ln \frac{0}{q} := 0$ for any $q \geq 0$, and $p \ln \frac{p}{0} := \infty$ for $p > 0$.
% \end{definition}
%
% \begin{definition}[Joint and conditional KL divergence]\label{def:joint_kl}
% Let $P$ and $Q$ be joint probability distributions over random variables $X, Y$ taking values in $\mathcal{X}, \mathcal{Y}$. The \textbf{joint} and \textbf{conditional} KL divergences are
% \[
% D_{\text{KL}}(P(X,Y) \,\|\, Q(X,Y)) := \sum_{x, y} P(x,y) \ln \frac{P(x,y)}{Q(x,y)},
% \qquad
% D_{\text{KL}}(P(Y \mid X) \,\|\, Q(Y \mid X)) := \sum_{x} P(x) \sum_{y} P(y \mid x) \ln \frac{P(y \mid x)}{Q(y \mid x)},
% \]
% with the same conventions as \cref{def:kl}.
% \end{definition}

The KL divergence between two joint distributions over a length-$n$ history splits, by an inductive application of the chain rule for probabilities, into a sum of per-step conditional KLs. This telescoping identity lets us trade a single joint KL over an entire history for a sum of per-step KLs (and vice versa): exactly the bridge we will need in \cref{prob:bound}.

We introduce two pieces of notation that will be used throughout the rest of the worksheet:

\begin{definition}[Cumulative and per-step KL]\label{def:Dn}
For environments / predictors $\nu, \rho$, horizon $n \geq 1$, and $1 \leq t \leq n$, the \textbf{cumulative joint KL} and the \textbf{expected per-step KL at step $t$} are
\begin{align*}
  D_n(\nu \parallel \rho) ~&:=~ \sum_{x_{1:n} \in \mathbb{B}^n}
                       \nu(x_{1:n}) \,\ln \frac{\nu(x_{1:n})}{\rho(x_{1:n})}, \\[6pt]
  d_t(\nu \parallel \rho) ~&:=~ \sum_{x_{<t} \in \mathbb{B}^{t-1}}
                       \nu(x_{<t}) \sum_{x_t \in \mathbb{B}}
                       \nu(x_t \mid x_{<t}) \,\ln \frac{\nu(x_t \mid x_{<t})}{\rho(x_t \mid x_{<t})},
\end{align*}
with the conventions $0 \ln \tfrac{0}{q} := 0$ for $q \geq 0$ and $p \ln \tfrac{p}{0} := \infty$ for $p > 0$. 
%The first argument always names the distribution the expectation is taken over (we do \emph{not} adopt Hutter's convention of an implicit ``$\mathbb{E}$ always over $\mu$,'' since the Pareto sections need $\nu$ to vary across the model class). 
%\Cref{prob:telescope_ex} below establishes that the two are related by $D_n(\nu \parallel \rho) = \sum_{t=1}^n d_t(\nu \parallel \rho)$.
We write $D_\infty(\nu \parallel \rho) := \lim_{t \to \infty} D_t(\nu \parallel \rho)$
and $D^\mu_\infty := D_\infty(\mu \parallel \xi)$.
\end{definition}

\begin{theorem}[KL divergence is non-negative]\label{thm:kl_nonneg}
For any environments $\nu, \rho$ and any horizon $n \geq 1$ and step $1 \leq t \leq n$:
\begin{description}\itemsep=2pt
  \item[\textnormal{(a)}] $D_n(\nu \parallel \rho) \geq 0$, with equality iff $\nu(x_{1:n}) = \rho(x_{1:n})$ for every $x_{1:n} \in \mathbb{B}^n$ with $\nu(x_{1:n}) > 0$.
  \item[\textnormal{(b)}] $d_t(\nu \parallel \rho) \geq 0$, with equality iff $\nu(x_t \mid x_{<t}) = \rho(x_t \mid x_{<t})$ for every $x_{<t}$ with $\nu(x_{<t}) > 0$ and every $x_t \in \mathbb{B}$.
\end{description}
Proof in \cref{app:kl_nonneg}.
\end{theorem}

\begin{exercise}[Telescoping KL]
\label{prob:telescope_ex}
\difficulty{10}



Show by induction on $n$ that for all environments $\nu, \rho$,
\[
D_n(\nu \parallel \rho) ~=~ \sum_{t=1}^{n} d_t(\nu \parallel \rho),
\]
where on the right-hand side, the $t=1$ summand uses the conventions $\nu(x_1 \mid x_{<1}) := \nu(x_1)$, $\rho(x_1 \mid x_{<1}) := \rho(x_1)$.

\begin{hint}
For the induction step, factor $\nu(x_{1:n}) = \nu(x_{<n})\,\nu(x_n \mid x_{<n})$ (and likewise for $\rho$) inside the log; the log splits additively into two pieces matching $D_{n-1}(\nu \parallel \rho)$ and $d_n(\nu \parallel \rho)$.
\end{hint}
\end{exercise}

\begin{solution}[prob:telescope_ex]
By induction on $n$.

\textbf{Base case ($n = 1$):}
\[
  D_1(\nu \parallel \rho) ~=~ \sum_{x_1} \nu(x_1) \ln \frac{\nu(x_1)}{\rho(x_1)} ~=~ d_1(\nu \parallel \rho),
\]
using $\nu(x_1 \mid x_{<1}) := \nu(x_1)$ and likewise for $\rho$.

\textbf{Induction step.} Apply the chain rule for probabilities (\cref{def:environment}) inside the log: $\nu(x_{1:n}) = \nu(x_{<n})\,\nu(x_n \mid x_{<n})$, and similarly for $\rho$. The log splits additively:
\begin{align*}
  & D_n(\nu \parallel \rho) \\
  ~&=~ \sum_{x_{1:n} \in \mathbb{B}^n} \nu(x_{1:n}) \,\ln \frac{\nu(x_{1:n})}{\rho(x_{1:n})} \\
  ~&=~ \sum_{x_{1:n}} \nu(x_n \mid x_{<n}) \nu(x_{<n}) \,\ln \frac{\nu(x_n \mid x_{<n})\,\nu(x_{<n})}
  {\rho(x_n \mid x_{<n})\,\rho(x_{<n})} \\
  ~&=~ \sum_{x_{1:n}} \nu(x_n \mid x_{<n}) \nu(x_{<n}) 
  \left[\ln \frac{\nu(x_n \mid x_{<n})}{\rho(x_n \mid x_{<n})} ~+~ \ln \frac{\nu(x_{<n})}{\rho(x_{<n})}\right]. \\
  ~&=~ \underbrace{\sum_{x_{<n}} \nu(x_{<n}) \sum_{x_n} \nu(x_n \mid x_{<n})  \ln \frac{\nu(x_n \mid x_{<n})}{\rho(x_n \mid x_{<n})} }_{d_n(\nu \parallel \rho)} \\
  &\qquad ~+~
  \sum_{x_{1:n}} \nu(x_n \mid x_{<n}) \nu(x_{<n}) \ln \frac{\nu(x_{<n})}{\rho(x_{<n})} \\
  ~&=~ d_n(\nu \parallel \rho) ~+~
  \underbrace{\sum_{x_{<n}}  \nu(x_{<n}) \ln \frac{\nu(x_{<n})}{\rho(x_{<n})}}_{D_{n-1}(\nu \parallel \rho)} \cancel{\sum_{x_n} \nu(x_n \mid x_{<n})} \\
  ~&=~ d_n(\nu \parallel \rho) + D_{n-1}(\nu \parallel \rho) 
  ~=~ d_n(\nu \parallel \rho) + \sum_{t=1}^{n-1} d_t(\nu \parallel \rho)  = \sum_{t=1}^{n} d_t(\nu \parallel \rho)
\end{align*}
% For the first piece, the log depends only on $x_{<n}$, so marginalizing $x_n$ collapses the outer sum to a sum over $x_{<n}$ weighted by $\nu(x_{<n})$ --- exactly $D_{n-1}(\nu \parallel \rho)$. For the second piece, factor $\nu(x_{1:n}) = \nu(x_{<n})\,\nu(x_n \mid x_{<n})$ to match the definition of $d_n(\nu \parallel \rho)$. Combining with the induction hypothesis $D_{n-1}(\nu \parallel \rho) = \sum_{t=1}^{n-1} d_t(\nu \parallel \rho)$:
% \[
%   D_n(\nu \parallel \rho) ~=~ D_{n-1}(\nu \parallel \rho) ~+~ d_n(\nu \parallel \rho) ~=~ \sum_{t=1}^{n} d_t(\nu \parallel \rho). \qedhere
% \]
\end{solution}

%===================================


% \[
% \int_p^q \frac{p-x}{x(1-x)}\, dx \;\leq\; 4\int_p^q (p-x)\, dx
% \;\;\Longrightarrow\;\;
% \int_q^p \frac{p-x}{x(1-x)}\, dx \;\geq\; 4\int_q^p (p-x)\, dx \;=\; 2(p-q)^2.
% \]
% Either way, $f_p(p) - f_p(q) \geq 2(p-q)^2$. \qedhere

% \medskip
% \emph{Remark (one-shot version).} Equivalently, note $f_p'(x) - 4(p-x) = (p-x)\big[\tfrac{1}{x(1-x)} - 4\big]$ has the same sign as $p-x$, since the bracket is $\geq 0$ on $(0,1)$. Hence the integrand is $\geq 0$ on $[q,p]$ when $q \leq p$ and $\leq 0$ on $[p,q]$ when $q > p$, so $\int_q^p \big[f_p'(x) - 4(p-x)\big]\, dx \geq 0$ in both cases.

% \medskip
% \emph{Alternative (no case split).} The same FTC computation avoids the case analysis entirely via a linear change of variables along the segment from $q$ to $p$. Using $2(p-q)^2 = 4\int_q^p (p-x)\, dx$,
% \[
% f_p(p) - f_p(q) - 2(p-q)^2 = \int_q^p \big[\,f_p'(x) - 4(p-x)\,\big]\, dx.
% \]
% Substitute $x = q + t(p-q)$ for $t \in [0,1]$, so $dx = (p-q)\, dt$ and $p - x = (1-t)(p-q)$:
% \[
% \int_q^p \big[\,f_p'(x) - 4(p-x)\,\big]\, dx
% = \int_0^1 (1-t)\,(p-q)^2 \Big[\tfrac{1}{x(1-x)} - 4\Big]\, dt
% \;\geq\; 0,
% \]
% since $(1-t) \geq 0$, $(p-q)^2 \geq 0$, and $\tfrac{1}{x(1-x)} - 4 \geq 0$ because $x = q + t(p-q) \in [0,1]$ is a convex combination of $q, p \in [0,1]$, so $x(1-x) \leq \tfrac14$. The orientation issue disappears because the $(p-x)$ in the integrand and the $(p-q)$ from $dx$ combine into $(p-q)^2 \geq 0$.


%===================================
\section{Mixture dominance}\label{prob:dominance}

The mixture $\xi$ never assigns much less probability than any single environment
weighted by its prior. This single inequality is the engine behind the cumulative bound
of \cref{prob:bound} below: dividing through gives $\mu(x)/\xi(x) \leq 1/w_\mu$, which
is then fed into Pinsker and the chain rule in the proof of the main cumulative bound.

\begin{exercise}[Mixture dominance]
\label{prob:dominance_ex}
\difficulty{05}



Show that for every $\nu \in \Mc$ and every $x \in \mathbb{B}^*$,
\[
  \xi(x) ~\geq~ w_\nu \cdot \nu(x).
\]
\end{exercise}

\begin{solution}[prob:dominance_ex]
By \cref{def:mixture}, $\xi(x) = \sum_{\nu' \in \Mc} w_{\nu'} \,\nu'(x)$.
Every term in this sum is non-negative, so dropping all terms except the one with
$\nu' = \nu$ can only decrease the sum:
\[
  \xi(x) ~=~ \sum_{\nu' \in \Mc} w_{\nu'}\, \nu'(x)
  ~\geq~ w_\nu\, \nu(x). \qedhere
\]
\end{solution}

\begin{remark}[Solomonoff specialization]
Under the Solomonoff prior, $w_\nu = 2^{-K(\nu)}$, so the bound becomes
\[
  \xi_U(x) ~\geq~ 2^{-K(\nu)} \cdot \nu(x) \qquad \text{for every computable }\nu.
\]
This is what makes $\xi_U$ \emph{universal}: a single predictor dominates the
entire computable model class up to a factor that depends only on the model's
description length.
\end{remark}

%===================================
\section{Cumulative prediction error bound (main result)}\label{prob:bound}


The next ingredient relates squared prediction error to KL divergence~\cite[Lemma~11.6.1]{Cover:06}. 
We require the following inequality:

\begin{theorem}[Pinsker's inequality]\label{thm:pinsker}
Let $P$ and $Q$ be probability distributions over $\mathbb{B} = \{\texttt{0}, \texttt{1}\}$, i.e., $q_0, q_1 \geq 0$
with $q_0 + q_1 = 1$. Write $p_x = P(x)$ and $q_x = Q(x)$. Then
\[
    \sum_{x \in \mathbb{B}} \big( q_x - p_x \big)^2 \;\leq\; \sum_{x \in \mathbb{B}} p_x \ln \frac{p_x}{q_x}.
    \tag{$\star$}
\]

\end{theorem}
The proof of this inequality is mostly tedious algebra but is included in \cref{app:pinsker}. 

% Let $\mu$ be any environment, and assume it generates our actual universe.

\begin{definition}[Cumulative expected and per-step squared prediction error]\label{def:S_infty}
For environments / predictors $\nu, \rho$ and horizon $n \geq 1$, the \textbf{cumulative expected squared error} and the \textbf{per-step squared error at step $t$} are
\begin{align*}
\Sn{\nu}{\rho} ~&:=~ \sum_{t=1}^n s_t(\nu \parallel \rho) \\
s_t(\nu \parallel \rho) ~&:=~ \sum_{x_{<t} \in \mathbb{B}^{t-1}} \nu(x_{<t}) \sum_{x_t \in \mathbb{B}} \bigl(\nu(x_t \mid x_{<t}) - \rho(x_t \mid x_{<t})\bigr)^2
\end{align*}
Write $S_{\infty}(\nu \parallel \rho) := \lim_{n \to \infty} \Sn{\nu}{\rho}$
and $S_\infty^\mu := S_{\infty}(\mu \parallel \xi)$.
\end{definition}

\begin{exercise}[Cumulative prediction bound]
\label{prob:bound_ex}
\difficulty{15}



Show that
\[
S_{\infty}^{\mu} \leq - \ln w_\mu.
\]

\begin{hint}
You need \cref{thm:pinsker}, \cref{prob:telescope_ex} and \cref{prob:dominance_ex}.
\end{hint}
\end{exercise}

\begin{solution}[prob:bound_ex]
The proof is a three-step chain:
\[
  S^\mu_\infty
  ~\stackrel{(1)}{\leq}~
   D^\mu_\infty
  ~\stackrel{(2)}{\leq}~
  \lim_{n \to \infty} \sum_{x_{1:n}} \mu(x_{1:n}) \,\ln \tfrac{1}{w_\mu}
  ~\stackrel{(3)}{=}~
  -\ln w_\mu.
\]

\medskip
\textit{Step (1): Pinsker, pointwise.}
For each $t$ and $x_{<t}$, both $\mu(\cdot \mid x_{<t})$ and $\xi(\cdot \mid x_{<t})$ are probability distributions on $\mathbb{B}$ ($\xi$ by \cref{prob:one_step}), so \cref{thm:pinsker} gives
\[
\sum_{x_t} \bigl(\xi(x_t \mid x_{<t}) - \mu(x_t \mid x_{<t})\bigr)^2 \leq \sum_{x_t} \mu(x_t \mid x_{<t}) \,\ln \tfrac{\mu(x_t \mid x_{<t})}{\xi(x_t \mid x_{<t})}.
\]
Multiplying by $\mu(x_{<t}) \geq 0$ and summing over $x_{<t}$ gives 
$s_t(\mu \parallel \xi) \leq d_t(\mu \parallel \xi)$, so $S_t(\mu \parallel \xi) \leq \sum_{i=1}^t d_i(\mu \parallel \xi) = D_t(\mu \parallel \xi)$ (\cref{prob:telescope_ex}), 
from which we take $t \to \infty$.

% \medskip
% \textit{Step (2): Telescoping.}
% By \cref{prob:telescope_ex}, $D_n(\mu \parallel \xi) = \sum_{t=1}^n d_t(\mu \parallel \xi)$. Taking $n \to \infty$ gives the equality.

\medskip
\textit{Step (2): Mixture dominance.}
By \cref{prob:dominance_ex}, $\xi(x_{1:n}) \geq w_\mu\, \mu(x_{1:n})$, so $\mu(x_{1:n}) / \xi(x_{1:n}) \leq 1/w_\mu$, hence $\ln(\mu/\xi) \leq \ln(1/w_\mu)$ pointwise. Substituting into the definition $D_n(\mu \parallel \xi) = \sum_{x_{1:n}} \mu(x_{1:n}) \ln(\mu/\xi)$ gives the bound.

\medskip
\textit{Step (3): Total mass.}
Pull the constant out and use $\sum_{x_{1:n}} \mu(x_{1:n}) = 1$ ($\mu$ is an environment). \qedhere
\end{solution}

%===================================

\begin{exercise}[Explicit complexity bound]
\label{prob:bound_explicit_ex}
\difficulty{05}



Specialize \cref{prob:bound} to the \textbf{Solomonoff prior} $w_\nu = 2^{-K(\nu)}$ (where $K(\nu)$ is the length of the shortest program computing $\nu$) to show
\[
S_{\infty}^{\mu} \leq K(\mu) \ln 2.
\]
\end{exercise}

\begin{solution}[prob:bound_explicit_ex]
By \cref{prob:bound}, for \emph{any} prior, $S_{\infty}^{\mu} \leq -\ln w_\mu$. The Solomonoff prior assigns $\mu$ the weight $w_\mu = 2^{-K(\mu)}$, so
\[
S_{\infty}^{\mu} ~\leq~ -\ln w_\mu ~=~ -\ln 2^{-K(\mu)} ~=~ K(\mu) \ln 2. \qedhere
\]
\end{solution}

\begin{exercise}[Per-step error]
\label{prob:per-step_error}
\difficulty{05}



Hence prove that the per-step expected prediction error
\[
s_t := \sum_{x_{<t} \in \mathbb{B}^{t-1}} \mu(x_{<t}) \sum_{x_t \in \mathbb{B}} \big( \xi(x_t \mid x_{<t}) - \mu(x_t \mid x_{<t}) \big)^2
\]
converges to zero as $t \to \infty$.
\end{exercise}

\begin{solution}[prob:per-step_error]
    Follows immediately as $\sum_{t=1}^\infty s_t \leq -\ln w_\mu < \infty$ implies $s_t \to 0$.
\end{solution}

% %===================================
% \section{From cumulative bound to a step-count}\label{prob:error_count}

% \Cref{prob:bound} gives an \emph{asymptotic} guarantee: the total
% expected squared prediction error is finite. We can turn this directly into a
% \emph{finite} statement: ``after at most this many steps, the mixture has
% predicted to within $\varepsilon$ on average.''

% \mysubsection{1}{[05] (Cumulative bound to step-count)\label{prob:error_count_ex}}

% Define the per-step expected squared error
% \[
%   d_t ~:=~ \sum_{x_{<t} \in \mathbb{B}^*} \mu(x_{<t})
%     \sum_{x_t \in \mathbb{B}} \big(\xi(x_t \mid x_{<t}) - \mu(x_t \mid x_{<t})\big)^2,
%   \qquad \text{so } S^\mu_\infty = \sum_{t=1}^\infty d_t.
% \]
% Show that for every $\varepsilon > 0$,
% \[
%   \#\big\{\, t \geq 1 \;:\; d_t > \varepsilon \,\big\}
%   ~<~ \frac{-\ln w_\mu}{\varepsilon}.
% \]

% \textit{Hint.} If $N$ such steps exist, what does that say about
% $N\varepsilon$ compared to $\sum_t d_t$?

% \begin{solution}
% Let $N := \#\{t : d_t > \varepsilon\}$. Each bad step contributes more than
% $\varepsilon$ to the cumulative sum, while all other steps contribute
% non-negatively, so
% \[
%   N \cdot \varepsilon ~<~ \sum_{t \, : \, d_t > \varepsilon} d_t
%   ~\leq~ \sum_{t=1}^\infty d_t ~=~ S^\mu_\infty ~\leq~ -\ln w_\mu,
% \]
% using \cref{prob:bound} at the last step. Dividing by $\varepsilon$ gives
% $N < -\ln w_\mu / \varepsilon$. \qedhere
% \end{solution}

% \begin{remark}[Solomonoff specialization]
% Under the Solomonoff prior $w_\mu = 2^{-K(\mu)}$, so $-\ln w_\mu = K(\mu)\ln 2$
% and the bound becomes
% \[
%   \#\big\{\, t \geq 1 \;:\; d_t > \varepsilon \,\big\} ~<~ \frac{K(\mu)\ln 2}{\varepsilon}.
% \]
% \end{remark}

% \begin{remark}
% Strictly fewer than $-\ln w_\mu / \varepsilon$ steps can ever see per-step
% expected error above $\varepsilon$; every other step is at most $\varepsilon$-bad
% on average. The bound is independent of the time-step $t$ and of any structure of
% the sequence beyond the prior weight of the true environment.
% \end{remark}

%===================================
\section{Bounded number of prediction mistakes}\label{prob:mistakes}

So far we have bounded \emph{squared error}. For a deterministic environment
(i.e.\ a single infinite binary sequence $x_1^*, x_2^*, \dots$ in $\Mc$),
one can ask the sharper $0/1$-loss question: how often does the mixture predict
the \emph{wrong bit}? The cumulative bound is strong enough to give an answer.


Let $\mu \in \Mc$ be a deterministic environment, so $\mu(x_t \mid x_{<t}) \in
\{0,1\}$ on every history; write $x_t^*$ for the unique bit on which
$\mu(\cdot \mid x^*_{<t})$ puts all its mass. Define the \textbf{threshold
predictor} associated with $\xi$: at each step, predict the more-likely bit,
\[
  \hat{x}_t ~:=~ \arg\max_{x \in \mathbb{B}} \xi(x \mid x_{<t})
  \qquad \text{(break ties arbitrarily).}
\]
We say the mixture makes a \textbf{mistake} at time $t$ if $\hat{x}_t \neq x_t^*$.

\begin{exercise}[Bounded prediction mistakes]
\label{prob:mistakes_ex}
\difficulty{15}



Show that the number of mistakes made by the threshold predictor on the
$\mu$-trajectory is at most
\[
  \#\{\, t \geq 1 \;:\; \hat{x}_t \neq x_t^* \,\}
  ~\leq~ -2 \ln w_\mu.
\]

\begin{hint}
At a mistake step, $\xi(x_t^* \mid x_{<t}) \leq \tfrac12$ (why?).
The per-step squared error along the deterministic $\mu$-trajectory
simplifies dramatically; bound it from below, then sum.
\end{hint}
\end{exercise}

\begin{solution}[prob:mistakes_ex]
\textbf{Step 1: per-step squared error along the $\mu$-trajectory.}
Because $\mu$ is deterministic, the outer sum over $x_{<t}$ in $d_t$ collapses
onto the unique trajectory $x_{<t}^* := x_1^* \cdots x_{t-1}^*$, with weight
$\mu(x_{<t}^*) = 1$. Write $\xi_t := \xi(x_t^* \mid x_{<t}^*) \in [0,1]$ for the
mass the mixture places on the \emph{correct} bit. Since $\mu$ assigns mass $1$
to $x_t^*$ and $0$ to the other bit,
\begin{align*}
  s_t ~&=~ \sum_{x_t \in \mathbb{B}} \big(\xi(x_t \mid x_{<t}^*) - \mu(x_t \mid x_{<t}^*)\big)^2 \\
  ~&=~ (\xi_t - 1)^2 + (1 - \xi_t - 0)^2 ~=~ 2(1 - \xi_t)^2.
\end{align*}

\textbf{Step 2: a mistake step contributes at least $\tfrac12$.}
A mistake at time $t$ means $\hat{x}_t \neq x_t^*$, i.e.\ the threshold predictor
chose the \emph{other} bit. This requires
$\xi(\hat{x}_t \mid x_{<t}^*) \geq \xi(x_t^* \mid x_{<t}^*)$, i.e.\
$1 - \xi_t \geq \xi_t$, i.e.\ $\xi_t \leq \tfrac12$. So at every mistake step,
\[
  s_t ~=~ 2(1 - \xi_t)^2 ~\geq~ 2 \cdot \tfrac{1}{4} ~=~ \tfrac12.
\]

\textbf{Step 3: sum and apply the cumulative bound.}
Let $N$ be the total number of mistakes. Summing $d_t$ over mistake steps
(non-negative contributions from non-mistake steps only help):
\[
  \tfrac12 \cdot N
  ~\leq~ \sum_{t \,:\, \hat{x}_t \neq x_t^*} s_t
  ~\leq~ \sum_{t=1}^\infty s_t ~=~ S^\mu_\infty
  ~\leq~ -\ln w_\mu,
\]
where the last step is \cref{prob:bound}. Rearranging, $N \leq -2 \ln w_\mu$.
\qedhere
\end{solution}

\begin{remark}[Solomonoff specialization]
Under the Solomonoff prior $w_\mu = 2^{-K(\mu)}$, so $-\ln w_\mu = K(\mu)\ln 2$
and the bound becomes
\[
  \#\{\, t \geq 1 \;:\; \hat{x}_t \neq x_t^* \,\} ~\leq~ 2 K(\mu) \ln 2.
\]
With $\ln 2 \approx 0.693$ the constant is just under $1.4$, so if the true 
environment $\mu$ is computable, and can be described by a program at most $k$ bits long, 
then the predictor $\xi$ 
will make at worst $\approx 1.4 k$ mistakes over predicting the entire sequence $x^*_{1:\infty}$.
Note that the bound is purely existence-style: it does not say \emph{when}
the mistakes happen. They could all occur at the start, or be arbitrarily far into the future.
\end{remark}

%===================================
\section{Pareto optimality under KL loss}\label{prob:kl}

The final two sections establish that the Bayesian mixture $\xi$ is
\emph{uniquely Pareto-optimal} under both KL and squared prediction loss:
no other predictor can do at least as well as $\xi$ on every environment
$\nu \in \Mc$ without coinciding with $\xi$ entirely. 
% Hutter~\cite[\S3.6]{Hutter:04uaibook} states a weaker (non-uniqueness) version of this result for several losses including the two we treat here; the Pythagorean-identity proofs we use give the stronger uniqueness statement essentially for free.

% A \textbf{predictor} (in this part) is any function $\rho$ assigning to each
% history $x_{<t}$ a distribution $\rho(\cdot \mid x_{<t}) \in \Delta\mathbb{B}$;
% by chain rule, $\rho(x_{1:n}) := \prod_t \rho(x_t \mid x_{<t})$. The mixture
% $\xi$ is one such predictor, and so is any $\nu \in \Mc$.

\begin{definition}[Pareto domination]\label{def:pareto}
For some measure of \textbf{loss} $\mathcal F(\nu, \rho)$ that takes an environment $\nu$ and a
predictor $\rho$, a predictor $\rho$ \textbf{weakly Pareto-dominates} $\xi$  with respect to 
loss $\mathcal{F}$ and class $\mathcal{M}$ if
\[
  \mathcal F(\nu, \rho) ~\leq~ \mathcal F(\nu, \xi)
  \qquad \text{for every } \nu \in \Mc.
\]
$\xi$ is \textbf{Pareto-optimal} (w.r.t.\ $\mathcal F$) if the only predictor
that weakly dominates it is $\xi$ itself.
%\footnote{This is stronger than the
%usual ``no strict domination'' formulation: weak domination already forces
%$\rho = \xi$, so a fortiori there is no $\rho$ that is strictly better on some
%$\nu$ without being worse on another.} 
% The quantifier ranges over the whole
% class --- there is no single ``true environment.'' Each $\nu$ takes its turn
% as the candidate environment; Pareto-optimality says no $\rho$ beats $\xi$
% uniformly across the class.

We specialize $\mathcal F$ to $D_n(\nu \parallel \rho)$ in \cref{prob:kl} 
and to $\Sn{\nu}{\rho}$ in \cref{prob:sq}.
\end{definition}


\begin{exercise}[KL Pythagorean identity]
\label{prob:kl_pythag}
\difficulty{10}



For \emph{any} predictor $\rho$ with $\rho(x_{1:n}) > 0$ on every $x_{1:n} \in \mathbb{B}^n$, show that
\[
  D_n(\xi \parallel \rho) ~+~ \sum_{\nu \in \Mc} w_\nu \, D_n(\nu \parallel \xi)
  ~=~ \sum_{\nu \in \Mc} w_\nu \, D_n(\nu \parallel \rho).
\]


\begin{hint}
Take the difference between the two summation terms.
\end{hint}
\end{exercise}

\begin{solution}[prob:kl_pythag]
Take the difference $\sum_\nu w_\nu D_n(\nu \parallel \rho) - \sum_\nu w_\nu D_n(\nu \parallel \xi)$.
Inside each $D_n$ the $\nu(x_{1:n}) \ln \nu(x_{1:n})$ pieces cancel, leaving
\begin{align*}
&\sum_\nu w_\nu D_n(\nu \parallel \rho) - \sum_\nu w_\nu D_n(\nu \parallel \xi) \\
&\quad= \sum_\nu w_\nu \sum_{x_{1:n}}
   \nu(x_{1:n})\,\Bigl[\ln\tfrac{\nu(x_{1:n})}{\rho(x_{1:n})}
                  - \ln\tfrac{\nu(x_{1:n})}{\xi(x_{1:n})}\Bigr] \\
&\quad= \sum_\nu w_\nu \sum_{x_{1:n}}
   \nu(x_{1:n})\,\ln \tfrac{\xi(x_{1:n})}{\rho(x_{1:n})}.
\end{align*}
The factor $\ln(\xi/\rho)$ does not depend on $\nu$, so we swap the order of
summation and pull it out, then use $\sum_\nu w_\nu \nu(x_{1:n}) = \xi(x_{1:n})$:
\begin{align*}
\sum_\nu w_\nu \sum_{x_{1:n}}
   \nu(x_{1:n})\,\ln \tfrac{\xi(x_{1:n})}{\rho(x_{1:n})}
&= \sum_{x_{1:n}} \ln \tfrac{\xi(x_{1:n})}{\rho(x_{1:n})}
   \underbrace{\sum_\nu w_\nu \nu(x_{1:n})}_{=\, \xi(x_{1:n})} \\
&= \sum_{x_{1:n}} \xi(x_{1:n}) \,\ln \tfrac{\xi(x_{1:n})}{\rho(x_{1:n})}
~=~ D_n(\xi \parallel \rho).
\end{align*}
Rearranging gives the claimed identity. \qedhere
\end{solution}

\begin{exercise}[KL Pareto-optimality]
\label{prob:kl_pareto}
\difficulty{05}



Show that $\xi$ is Pareto-optimal under $D_n(\nu \parallel \rho)$ in the sense of \cref{def:pareto}: if $\rho$ is any predictor with
\[
  D_n(\nu \parallel \rho) ~\leq~ D_n(\nu \parallel \xi) \qquad \text{for every } \nu \in \Mc,
\]
then $\rho = \xi$.

\begin{hint}
Multiply the assumed inequality by $w_\nu \geq 0$ and sum over
$\nu \in \Mc$; then apply the KL Pythagorean identity to recognize the extra
non-negative term, which must vanish.
\end{hint}
\end{exercise}

\begin{solution}[prob:kl_pareto]
Multiply the assumed inequality by $w_\nu \geq 0$ and sum over $\nu \in \Mc$:
\begin{equation}\label{eq:kl_weighted_dom}
  \sum_\nu w_\nu\, D_n(\nu \parallel \rho) ~\leq~ \sum_\nu w_\nu\, D_n(\nu \parallel \xi).
\end{equation}
The KL Pythagorean identity rewrites the LHS as
\[
  D_n(\xi \parallel \rho) \;+\; \sum_\nu w_\nu\, D_n(\nu \parallel \xi),
\]
so \cref{eq:kl_weighted_dom} forces $D_n(\xi \parallel \rho) \leq 0$. KL is non-negative, hence $D_n(\xi \parallel \rho) = 0$, which forces $\rho(x_{1:n}) = \xi(x_{1:n})$ for every $x_{1:n} \in \mathbb{B}^n$. By chain rule, the conditionals agree at every history $x_{<t}$ with $\xi(x_{<t}) > 0$. \qedhere
\end{solution}


%===================================
\section{Misspecified models: when $\mu \notin \Mc$}\label{prob:misspec}

The cumulative bound \cref{prob:bound} assumed the true environment $\mu$ lives
in the model class $\Mc$. What happens when it does not? We now show the
bound \emph{degrades gracefully}: it splits cleanly into a complexity term
(``cost of not knowing which model is best''), exactly as before, plus a
linear-in-time approximation term (``cost of no model being right''). See \cite[\S3.2.8]{Hutter:04uaibook} for the original treatment.

Throughout this section, we no longer assume $\mu \in \Mc$. 
We will state the bound in terms of an arbitrary
$\hat\mu \in \Mc$, so taking the infimum over $\hat\mu$ on the
right-hand side gives the tightest version by using the ``closest''
approximation of $\hat{\mu} \in \Mc$ to $\mu$.
\begin{exercise}[Misspecified KL bound]
\label{prob:misspec_ex}
\difficulty{10}



Show that for all $\hat\mu \in \Mc$,
\[
  D_n(\mu \parallel \xi) ~\leq~ -\ln w_{\hat\mu} ~+~ D_n(\mu \parallel \hat\mu).
\]

\begin{hint}
Use \cref{prob:dominance_ex}.
\end{hint}
\end{exercise}

\begin{solution}[prob:misspec_ex]
Mixture dominance applied to $\hat\mu \in \Mc$ (\cref{prob:dominance}) gives
$\xi(x_{1:n}) \geq w_{\hat\mu}\, \hat\mu(x_{1:n})$ for every $x_{1:n}$. Hence
$\mu(x_{1:n}) / \xi(x_{1:n}) \leq \mu(x_{1:n}) / (w_{\hat\mu}\, \hat\mu(x_{1:n}))$, and taking logs,
\[
  \ln \frac{\mu(x_{1:n})}{\xi(x_{1:n})}
  ~\leq~ -\ln w_{\hat\mu} \;+\; \ln \frac{\mu(x_{1:n})}{\hat\mu(x_{1:n})}.
\]
Now take the $\mu$-expectation. The LHS becomes $D_n(\mu \parallel \xi)$; the second RHS term becomes $D_n(\mu \parallel \hat\mu)$; the constant $-\ln w_{\hat\mu}$ survives because $\mu(X_{1:n})$ has total mass $1$. \qedhere
\end{solution}

\begin{remark}[Reading the two terms]
The right-hand side splits into two qualitatively different terms:
\begin{itemize}\itemsep=2pt
\item $-\ln w_{\hat\mu}$ is \emph{constant in $n$}: under the Solomonoff
prior, this is $K(\hat\mu)\ln 2$. The familiar ``complexity'' cost of search.
\item The approximation term $D_n(\mu \parallel \hat\mu) = \sum_{t=1}^n d_t(\mu \parallel \hat\mu)$ (by \cref{prob:telescope_ex}) is a sum of $n$ per-step KLs; in general it grows with $n$. It vanishes identically iff $\hat\mu$ matches $\mu$ on every $\mu$-reachable history, in particular when $\mu \in \Mc$ and we pick $\hat\mu = \mu$, recovering \cref{prob:bound}.
\end{itemize}
Combining with Pinsker (\cref{thm:pinsker}) gives the corresponding bound on
$S^\mu_\infty$, which is infinite in general but inherits the rate of $\hat\mu$.
\end{remark}

Details on how to define the best choice of $\hat{\mu}$ are in \cref{app:best_choice_mu}.

%===================================
\section{Pareto optimality under squared loss}\label{prob:sq}

The squared specialization $\mathcal F(\nu, \rho) = \Sn{\nu}{\rho}$ (\cref{def:S_infty}) lives one level down from KL: at the conditional distributions
$\nu(\cdot \mid x_{<t})$. The natural Pythagorean decomposition at a fixed
history uses \emph{posterior} weights $w(\nu \mid x_{<t})$: those are the
weights that make $\xi(x_t \mid x_{<t})$ the mean of the $\nu(x_t \mid x_{<t})$'s
(via the posterior-predictive form, \cref{def:mixture}). The Pareto-optimality
aggregation, however, uses prior weights $w_\nu$. Bridging the two takes one
extra Bayes-rule step.

\begin{exercise}[Squared Pythagorean identity]
\label{prob:sq_pythag}
\difficulty{10}



Fix any $t \in \{1, \dots, n\}$ and any history $x_{<t} \in \mathbb{B}^{t-1}$:
\emph{the history is arbitrary, not sampled from any environment}. For any
$x_t \in \mathbb{B}$ and any predictor $\rho$, show
\begin{align*}
  &\bigl(\xi(x_t \mid x_{<t}) - \rho(x_t \mid x_{<t})\bigr)^{2}
  ~+~
  \sum_{\nu \in \Mc} w(\nu \mid x_{<t})\,\bigl(\nu(x_t \mid x_{<t}) - \xi(x_t \mid x_{<t})\bigr)^{2} \\
  &=~
  \sum_{\nu \in \Mc} w(\nu \mid x_{<t})\,\bigl(\nu(x_t \mid x_{<t}) - \rho(x_t \mid x_{<t})\bigr)^{2}.
\end{align*}

\begin{hint}
Add and subtract $\xi(x_t \mid x_{<t})$ inside the squared term,
then expand. Use $\sum_\nu w(\nu \mid x_{<t}) = 1$ and the posterior-predictive
form of $\xi$ (\cref{def:mixture}), $\sum_\nu w(\nu \mid x_{<t})\,\nu(x_t \mid x_{<t})
= \xi(x_t \mid x_{<t})$, to kill the cross-term. The weighting must be
\emph{posterior} weights $w(\nu \mid x_{<t})$ (not prior weights $w_\nu$),
because $\xi(x_t \mid x_{<t})$ is the posterior-weighted mean of the
$\nu(x_t \mid x_{<t})$'s, not the prior one.
\end{hint}
\end{exercise}

\begin{solution}[prob:sq_pythag]
Throughout, the history $x_{<t}$ and symbol $x_t$ are fixed but arbitrary.
Abbreviate $\nu := \nu(x_t \mid x_{<t})$, $\xi := \xi(x_t \mid x_{<t})$,
$\rho := \rho(x_t \mid x_{<t})$, and $\tilde w_\nu := w(\nu \mid x_{<t})$ for
the duration of the calculation.

\medskip
\textit{Step 1: $\xi$ is the posterior-weighted mean of the $\nu$'s.}
By the posterior-predictive form of the mixture (\cref{def:mixture}) and the
posterior normalization $\sum_\nu \tilde w_\nu = 1$,
\[
  \sum_\nu \tilde w_\nu \, \nu ~=~ \xi.
\]

\medskip
\textit{Step 2: the cross-term vanishes.}
A direct consequence of Step 1, again using $\sum_\nu \tilde w_\nu = 1$:
\[
  \sum_\nu \tilde w_\nu \, (\nu - \xi)
  ~=~ \sum_\nu \tilde w_\nu \, \nu \;-\; \xi \sum_\nu \tilde w_\nu
  ~=~ \xi - \xi ~=~ 0.
\]

\medskip
\textit{Step 3: expand the square.}
Write $\nu - \rho = (\nu - \xi) + (\xi - \rho)$ and expand:
\begin{align*}
&\sum_\nu \tilde w_\nu (\nu - \rho)^2 \\
&\quad= \sum_\nu \tilde w_\nu\bigl[(\nu - \xi)^2 + 2(\nu - \xi)(\xi - \rho) + (\xi - \rho)^2\bigr] \\
&\quad= \sum_\nu \tilde w_\nu (\nu - \xi)^2
   \;+\; 2(\xi - \rho)\sum_\nu \tilde w_\nu (\nu - \xi)
   \;+\; (\xi - \rho)^2 \sum_\nu \tilde w_\nu.
\end{align*}
The middle sum is zero by Step 2; the trailing $\sum_\nu \tilde w_\nu = 1$. So
\[
  (\xi - \rho)^2 \;+\; \sum_\nu \tilde w_\nu (\nu - \xi)^2
  ~=~ \sum_\nu \tilde w_\nu (\nu - \rho)^2. \qedhere
\]
\end{solution}

\begin{exercise}[Squared Pareto-optimality]
\label{prob:sq_pareto}
\difficulty{05}



Show that $\xi$ is Pareto-optimal under $\Sn{\nu}{\rho}$ in the sense of \cref{def:pareto}: if $\rho$ is any predictor with
\[
  \Sn{\nu}{\rho} ~\leq~ \Sn{\nu}{\xi} \qquad \text{for every } \nu \in \Mc,
\]
then $\rho(x_t \mid x_{<t}) = \xi(x_t \mid x_{<t})$ at every history $x_{<t} \in
\mathbb{B}^{t-1}$ \emph{with $\xi(x_{<t}) > 0$} (equivalently, $\xi$-almost
surely) and every $x_t \in \mathbb{B}$. Histories that the mixture never reaches
($\xi(x_{<t}) = 0$, i.e., reached by no $\nu \in \Mc$ either) are invisible
to the loss and are not pinned down.

\begin{hint}
Multiply the assumed inequality by $w_\nu \geq 0$ and sum over
$\nu$; then for each fixed history $x_{<t}$, multiply the per-history
Pythagorean identity from above by $\xi(x_{<t})$ and use the Bayes identity
$w(\nu \mid x_{<t})\,\xi(x_{<t}) = w_\nu \,\nu(x_{<t})$ to bridge from posterior
weights (inside the identity) to prior weights (in the aggregated Pareto
inequality).
\end{hint}
\end{exercise}

\begin{solution}[prob:sq_pareto]
Multiply the assumed inequality by $w_\nu \geq 0$ and sum over $\nu \in \Mc$:
\begin{equation}\label{eq:sq_weighted_dom}
  \sum_\nu w_\nu\,\Sn{\nu}{\rho}
  ~\leq~
  \sum_\nu w_\nu\,\Sn{\nu}{\xi}.
\end{equation}

Now bridge to the per-history Pythagorean. At each \emph{fixed} history $x_{<t}
\in \mathbb{B}^{t-1}$ and symbol $x_t \in \mathbb{B}$, the identity says
\[
  (\xi - \rho)^2 \;+\; \sum_\nu w(\nu \mid x_{<t})\,(\nu - \xi)^2
  ~=~ \sum_\nu w(\nu \mid x_{<t})\,(\nu - \rho)^2,
\]
with the abbreviations $\nu := \nu(x_t \mid x_{<t})$, $\xi := \xi(x_t \mid x_{<t})$,
$\rho := \rho(x_t \mid x_{<t})$. Multiply by $\xi(x_{<t})$ and use Bayes,
$w(\nu \mid x_{<t})\,\xi(x_{<t}) = w_\nu \,\nu(x_{<t})$:
\[
  \xi(x_{<t})(\xi - \rho)^2
  \;+\; \sum_\nu w_\nu \,\nu(x_{<t})\,(\nu - \xi)^2
  ~=~ \sum_\nu w_\nu \,\nu(x_{<t})\,(\nu - \rho)^2.
\]
Sum over $t = 1, \dots, n$ and all histories $x_{<t} \in \mathbb{B}^{t-1}$ and
symbols $x_t \in \mathbb{B}$. The two $\sum_\nu w_\nu \,\nu(x_{<t})(\cdots)$
terms become exactly $\sum_\nu w_\nu \,\Sn{\nu}{\xi}$ and $\sum_\nu w_\nu \,
\Sn{\nu}{\rho}$ respectively, so
\[
  \|\xi - \rho\|_\xi^2
  \;+\; \sum_\nu w_\nu\,\Sn{\nu}{\xi}
  ~=~
  \sum_\nu w_\nu\,\Sn{\nu}{\rho},
\]
where
\[
  \|\xi - \rho\|_\xi^2 ~:=~
  \sum_{t=1}^n \sum_{x_{<t}} \xi(x_{<t}) \sum_{x_t}
    \bigl(\xi(x_t \mid x_{<t}) - \rho(x_t \mid x_{<t})\bigr)^2 ~\geq~ 0.
\]
Combining with \cref{eq:sq_weighted_dom},
\[
  \|\xi - \rho\|_\xi^2 \;+\; \sum_\nu w_\nu\,\Sn{\nu}{\xi}
  ~\leq~ \sum_\nu w_\nu\,\Sn{\nu}{\xi}
  \;\;\Longrightarrow\;\;
  \|\xi - \rho\|_\xi^2 ~\leq~ 0.
\]
Since $\|\xi - \rho\|_\xi^2$ is a sum of non-negative terms and is itself
$\leq 0$, every term vanishes:
$\xi(x_{<t})\bigl(\xi(x_t \mid x_{<t}) - \rho(x_t \mid x_{<t})\bigr)^2 = 0$
at every $(t, x_{<t}, x_t)$. So $\rho(x_t \mid x_{<t}) = \xi(x_t \mid x_{<t})$
at every history $x_{<t}$ with $\xi(x_{<t}) > 0$. \qedhere
\end{solution}

%===================================
\clearpage
\appendix
\renewcommand{\thesection}{\Alph{section}}
\crefname{section}{Appendix}{Appendices}
\Crefname{section}{Appendix}{Appendices}

\section{Proof of Pinsker's inequality}\label{app:pinsker}

We first prove the following lemma:

\begin{lemma}[Pinsker's binary inequality]\label{lem:pinsker_binary}
    For $0 \leq p \leq 1$ and $0 < q < 1$ we have
    \[
        2(p-q)^2 ~\leq~ p \ln \frac{p}{q} + (1-p) \ln \frac{1-p}{1-q}.
    \]
    \end{lemma}
    
    % \mysubsection{1}{[15] (Pinsker binary inequality)\label{prob:pinsker_ex}}
    
    % Prove \cref{lem:pinsker_binary}.
    
    % \textit{Hint:} Define a function $f(x)$ such that the right-hand side can be written as $f(p) - f(q)$ (try separating the LHS into terms with and without $q$). Then use the fundamental theorem of calculus, $f(p) - f(q) = \int_q^p f'(x) \, dx$, together with the inequality $x(1-x) \leq 1/4$ for $0 \leq x \leq 1$, to obtain the result.
    % Alternatively: Fix $0 \leq p \leq 1$ and let $\phi_p(q) := p \ln \frac{p}{q} + (1-p) \ln \frac{1-p}{1-q} - 2(p-q)$ and show for any 
    % $0 \leq p \leq 1$, $\phi_p(q)$ has a global minimum at $q = p$. 
    
    Define $f(x) := p \ln x + (1-p) \ln (1-x)$, so that
    \[
    p \ln \frac{p}{q} + (1-p) \ln \frac{1-p}{1-q} = f(p) - f(q)
    = \int_q^p f'(x)\, dx
    = \int_q^p \frac{p-x}{x(1-x)}\, dx,
    \]
    using $f_p'(x) = \frac{p}{x} - \frac{1-p}{1-x} = \frac{p-x}{x(1-x)}$. Recall also that $x(1-x) \leq \tfrac14$ on $[0,1]$, so $\frac{1}{x(1-x)} \geq 4$.
    
    \textbf{Case $q \leq p$.} For every $x \in [q,p]$ we have $x \leq p$, hence $p - x \geq 0$.
    \[
    \frac{p-x}{x(1-x)} \;\geq\; 4(p-x) \qquad \text{for } x \in [q,p].
    \]
    Integrating over $[q,p]$,
    \[
    \int_q^p \frac{p-x}{x(1-x)}\, dx \;\geq\; 4 \int_q^p (p-x)\, dx \;=\; 2(p-q)^2.
    \]
    
    \textbf{Case $q \geq p$.} For every $x \in [p,q]$ we have $x \geq p$, hence $x - p \geq 0$. 
    \[
    \frac{x-p}{x(1-x)} \;\geq\; 4(x-p) \qquad \text{for } x \in [p,q].
    \]
    Integrating over $[q,p]$,
    \begin{align*}
    \int_q^p \frac{p-x}{x(1-x)}\, dx
    &= \int_p^q \frac{x-p}{x(1-x)}\, dx \\
    &\geq 4 \int_p^q (x-p)\, dx
    = 4 \int_q^p (p-x)\, dx
    = 2(p-q)^2
    \end{align*}

\begin{proof}[Proof of \cref{thm:pinsker}]
    We now reduce \cref{thm:pinsker} to the binary inequality \cref{lem:pinsker_binary}.
    Recall $p_0 + p_1 = 1$ and $q_0 + q_1 = 1$. We handle the boundary case
$q_x = 0$ first, then reduce the interior case directly to
\cref{lem:pinsker_binary}.

\medskip
\textbf{Case 1: $q_x = 0$ for some $x \in \mathbb{B}$.}
If $p_x > 0$, then the right-hand side of $(\star)$ contains
$p_x \ln \tfrac{p_x}{0} = +\infty$, so $(\star)$ holds trivially.
If instead $p_x = 0$, that atom contributes $0$ to both sides
(using the convention $0 \ln \tfrac{0}{0} = 0$). The other atom $x'$ then
satisfies $p_{x'} = 1$ and $q_{x'} = 1 - q_x = 1$, so both sides vanish and
$(\star)$ holds.

\medskip
\textbf{Case 2: $q_0, q_1 > 0$.}
Since $q_1 = 1 - q_0$ and $p_1 = 1 - p_0$, the two sides of $(\star)$ become
\[
    \sum_{x \in \mathbb{B}} (q_x - p_x)^2
    = (q_0 - p_0)^2 + \big((1-q_0) - (1-p_0)\big)^2
    = 2(p_0 - q_0)^2,
\]
\[
    \sum_{x \in \mathbb{B}} p_x \ln \tfrac{p_x}{q_x}
    = p_0 \ln \tfrac{p_0}{q_0} + (1 - p_0) \ln \tfrac{1 - p_0}{1 - q_0}.
\]
So $(\star)$ is exactly
$2(p_0 - q_0)^2 \leq p_0 \ln \tfrac{p_0}{q_0} + (1 - p_0) \ln \tfrac{1 - p_0}{1 - q_0}$,
which is \cref{lem:pinsker_binary} with $p = p_0$ and $q = q_0$
(here $0 < q_0 < 1$, since $q_0, q_1 > 0$).
\end{proof}

%===================================
\section{Proof of KL non-negativity}\label{app:kl_nonneg}

We prove \cref{thm:kl_nonneg}. The core is the elementary inequality
\[
  \ln x ~\leq~ x - 1 \qquad \text{for all } x > 0, \quad \text{equality iff } x = 1.
\]
(Let $f(x) := (x - 1) - \ln x$; then $f'(x) = 1 - 1/x$, vanishing only at $x = 1$, with $f''(x) = 1/x^2 > 0$, so $f$ is strictly convex and attains its unique minimum $f(1) = 0$ at $x = 1$.)

\begin{lemma}[Gibbs' inequality]\label{lem:gibbs}
Let $P, Q$ be probability distributions over a finite set $\mathcal{X}$, with the conventions $0 \ln \tfrac{0}{q} := 0$ and $p \ln \tfrac{p}{0} := \infty$. Then
\[
  \sum_{x \in \mathcal{X}} P(x) \ln \frac{P(x)}{Q(x)} ~\geq~ 0,
\]
with equality iff $P(x) = Q(x)$ for every $x$ with $P(x) > 0$.
\end{lemma}

\begin{proof}
If $Q(x) = 0$ for some $x$ with $P(x) > 0$, the LHS is $+\infty$ and the inequality is strict. Otherwise restrict the sum to $\{x : P(x) > 0\}$ (terms with $P(x) = 0$ contribute $0$ by convention), where every ratio $Q(x)/P(x)$ is in $(0, \infty)$. Apply $\ln y \leq y - 1$ to $y = Q(x)/P(x)$:
\begin{align*}
  - \sum_x P(x) \ln \frac{P(x)}{Q(x)}
  ~=~ \sum_x P(x) \ln \frac{Q(x)}{P(x)}
  ~\leq~ \sum_x P(x) \left(\frac{Q(x)}{P(x)} - 1\right) \\
  ~=~ \sum_x Q(x) - \sum_x P(x) ~\leq~ 0,
\end{align*}
where the last step uses $\sum_x Q(x) \leq 1 = \sum_x P(x)$ (the sum over $\{x : P(x) > 0\}$ may miss some $Q$-mass). Negating gives the claim. Equality in $\ln y \leq y - 1$ requires $y = 1$ for every $x$ with $P(x) > 0$, i.e.\ $Q(x) = P(x)$ there; equality in $\sum Q \leq \sum P$ then forces $Q$ to have no mass outside $\mathrm{supp}\,P$. Both together give $P = Q$ on all of $\mathcal{X}$.
\end{proof}

\begin{proof}[Proof of \cref{thm:kl_nonneg}]
\textit{Part (a).} Apply \cref{lem:gibbs} with $\mathcal{X} = \mathbb{B}^n$, $P(x_{1:n}) = \nu(x_{1:n})$, $Q(x_{1:n}) = \rho(x_{1:n})$. Both are probability distributions on $\mathbb{B}^n$ since $\nu$ and $\rho$ are environments. The lemma gives $D_n(\nu \parallel \rho) \geq 0$ with the stated equality condition.

\textit{Part (b).} For each fixed history $x_{<t}$ with $\nu(x_{<t}) > 0$, apply \cref{lem:gibbs} with $\mathcal{X} = \mathbb{B}$, $P = \nu(\cdot \mid x_{<t})$, $Q = \rho(\cdot \mid x_{<t})$ (both are probability distributions on $\mathbb{B}$):
\[
  \sum_{x_t \in \mathbb{B}} \nu(x_t \mid x_{<t}) \ln \frac{\nu(x_t \mid x_{<t})}{\rho(x_t \mid x_{<t})} ~\geq~ 0,
\]
with equality iff $\nu(\cdot \mid x_{<t}) = \rho(\cdot \mid x_{<t})$. Multiply by $\nu(x_{<t}) \geq 0$ and sum over $x_{<t}$: this gives $d_t(\nu \parallel \rho) \geq 0$, as a sum of non-negative terms. Equality forces every $x_{<t}$ with $\nu(x_{<t}) > 0$ to contribute a vanishing inner sum, which is exactly the stated condition.
\end{proof}

%===================================

\section{Model Class $\Mc$}\label{app:model_class}

All results in this sheet hold for any countable class $\Mc$ and prior $w_\nu$. The canonical \textbf{Solomonoff} choice takes $\Mc$ to be the class of all computable environments and the \textbf{Solomonoff prior}
\[
w_\nu ~:=~ 2^{-K(\nu)},
\]
where $K(\nu)$ is the \emph{Kolmogorov complexity} of $\nu$: the length of the shortest program that computes $\nu$. The Bayesian mixture $\xi$ under this particular prior is the \emph{universal mixture}, written $\xi_U$. With this class the assumption $\mu \in \Mc$ reduces to ``the universe is computable'': as weak an assumption as one can make. For technical details on Kolmogorov complexity, see \cite[\S2.7]{Hutter:24uaibook2}.

The class of all computable measures is countable but not enumerable (by the halting problem, you can't decide which Turing machines define total functions), so one can't even formally write
$\sum_{\nu \in \Mc} \ldots$, as the decision problem of determining if $\nu \overset{?}{\in} \Mc$ is undecidable. A faithful construction of $\xi_U$ requires expanding $\Mc$ to all \emph{lower-semicomputable semimeasures} (Levin), at which point $\xi_U$ becomes a semimeasure rather than a measure: $\sum_x \xi_U(x) \leq 1$ instead of $=1$, conditionals may not sum to $1$, and every other proof in this sheet gains a side-case. This also gives the nice benefit that $\xi \in \mathcal{M}$.

We dodge all of this: $\Mc$ is just \emph{some} countable set of proper measures with $\mu \in \Mc$, and we assume $K(\nu)$ is defined for each $\nu \in \Mc$ when we need it. We never need $\xi$ itself to be in $\Mc$, nor do we need its computability, so we don't ask. See \cite[\S2.4.3]{Hutter:04uaibook} for the full Levin construction.

%===================================

\section{Best choice of $\hat{\mu} \in \Mc$.}\label{app:best_choice_mu}

The bound is valid for every $\hat\mu \in \Mc$. A tempting question: which
$\hat\mu$ gives the tightest bound? The answer is genuinely setting-dependent.

\textit{Finite horizon $n$.} The minimizer of the RHS is
\[
    \hat\mu_n ~:=~ \arg\min_{\nu \in \Mc}\Bigl\{ -\ln w_\nu \;+\; D_n(\mu \parallel \nu) \Bigr\}.
\]
This is well-defined (finitely many terms, attained when $\Mc$ is finite),
but the answer depends on $n$: for small $n$ the prior term dominates and a
high-prior coarse model wins; for large $n$ the fit term dominates and the best
per-step approximator wins. Both regimes are correct.

\textit{Asymptotic rate.} A horizon-free analogue minimizes the per-step KL
rate $$\lim_{n \to \infty} \frac{1}{n} \sum_{t \leq n} d_t(\mu \parallel \nu).$$ This is the right notion when one cares about the linear-growth slope of the bound. But it has two drawbacks: the limit need not exist for non-stationary $\mu$ (one can substitute $\liminf$ but at the cost of clarity), and it discards the prior $w_\nu$ entirely, so the minimizer is not the same as $\hat\mu_n$ at any finite $n$, since only the slope is minimized.

\textit{Stationary $\mu$.} The per-step KLs $d_t(\mu \parallel \nu)$ are constant in $t$, so all three notions essentially agree: $\hat\mu = \arg\min_{\nu \in \Mc} D_{\text{KL}}(\mu \,\|\, \nu)$, where the KL is taken under the (common) one-step conditional law. This is the case where ``$\hat\mu$ = projection of $\mu$ onto $\Mc$ in KL'' has unambiguous meaning.

\textit{Upshot.} The bound holds for all $\hat\mu \in \Mc$, and one is free
to take the infimum on the RHS over $\hat\mu$. Which $\hat\mu$ actually
attains that infimum is a separate question whose answer depends on horizon,
prior, and any structural assumptions on $\mu$.
%===================================

\bibliographystyle{alphaurl}
\bibliography{biblo}

\end{document}
